\documentclass[11pt]{article}
\usepackage[margin=1in]{geometry}
\usepackage{amsmath,amssymb,amsthm,mathtools}
\usepackage{graphicx}
\usepackage[hidelinks]{hyperref}

\newtheorem{theorem}{Theorem}
\newtheorem{lemma}[theorem]{Lemma}
\newtheorem{corollary}[theorem]{Corollary}
\newtheorem{remark}[theorem]{Remark}
\newcommand{\M}{\mathcal M}
\newcommand{\A}{\mathcal A}
\newcommand{\D}{\mathcal D}
\newcommand{\E}{\mathbb E}
\newcommand{\BMO}{\mathrm{BMO}}
\newcommand{\dd}{\,dm}

\title{A commuting-symbol characterization for\\
semicommutative martingale paraproducts}
\author{Automated mathematics run \texttt{fa\_banach\_001}, lane 8}
\date{19 July 2026}

\begin{document}
\maketitle

\paragraph{Status.}
Partial result, likely valid, subject to human review.  It settles a natural
commuting-symbol subcase of the open characterization problem, not the general
noncommutative problem.

\section{Source question}

Wei and Zhang ask how to characterize boundedness of an operator-valued
martingale paraproduct on semicommutative $L_2$ \cite[p.~2]{WeiZhang}:

\begin{quote}
``It still remains open how to characterize the boundedness of $\pi_b$ on
$L_2(L_\infty(\mathbb R)\overline\otimes\M)$.''
\end{quote}

The surrounding source passage, including the finite-dimensional logarithmic
obstruction and the definition of the symbol class, is reproduced below.

\begin{center}
\includegraphics[width=\textwidth]{figures/open_problem_crop.png}
\end{center}

Let $d\ge2$, let $\D$ be the $d$-adic filtration on $\mathbb R$, and write
$\E_k$ and $d_k=\E_k-\E_{k-1}$ for conditional expectation and martingale
difference.  If $(\M,\tau)$ is a semifinite von Neumann algebra, the
semicommutative paraproduct is initially defined on finite martingales by
\[
  \pi_b f=\sum_k d_kb\,\E_{k-1}f.
\]
We use the column martingale BMO seminorm
\begin{equation}\label{eq:bmoc}
 \|b\|_{\BMO_{c,d}}
 :=\sup_n\left\|\E_n\sum_{k\ge n}|d_kb|^2\right\|_{\M}^{1/2}.
\end{equation}
Changing $k\ge n$ to $k>n$ gives an equivalent seminorm, with constants
depending only on $d$; this elementary indexing point is included in the
proof below.

\section{Idea of the proof}

The obstruction in the general problem is noncommutation of the operator
coefficients.  If all values of $b$ lie in one abelian von Neumann algebra,
left multiplication by those values is a commuting normal family on
$L_2(\M)$.  The abelian spectral theorem simultaneously diagonalizes that
family.  Each finite truncation of $\pi_b$ then becomes a direct integral of
scalar paraproducts, tensored with identity operators on multiplicity spaces.

The operator norm of this direct integral is the essential supremum of its
scalar fiber norms.  Positivity and the same diagonalization show that
\eqref{eq:bmoc} is precisely the essential supremum of the scalar fiber BMO
seminorms.  The scalar Carleson embedding theorem compares the scalar
paraproduct and BMO norms.  Uniform bounds for all truncations give the full
operator.  Thus the ambient algebra and the input may remain noncommutative;
only the range of the symbol is required to commute.

\section{Formal result}

\begin{theorem}[Commuting-range characterization]\label{thm:main}
Let $(\M,\tau)$ be a semifinite von Neumann algebra, and let
$b:\mathbb R\to\M$ be locally Bochner integrable and adapted to the $d$-adic
filtration.  Suppose there is an abelian von Neumann subalgebra
$\A\subseteq\M$ such that $b(t)\in\A$ for almost every $t$.

Then $\pi_b$ extends boundedly to
\[
 L_2(L_\infty(\mathbb R)\overline\otimes\M)
 \cong L_2(\mathbb R;L_2(\M,\tau))
\]
if and only if $\|b\|_{\BMO_{c,d}}<\infty$.  Moreover, there are constants
$0<c_d\le C_d<\infty$, depending only on $d$, such that
\begin{equation}\label{eq:equiv}
 c_d\|b\|_{\BMO_{c,d}}
 \le \|\pi_b\|_{2\to2}
 \le C_d\|b\|_{\BMO_{c,d}}.
\end{equation}
\end{theorem}

\begin{corollary}
The conclusion holds whenever the operators in the essential range of $b$
are normal and commute pairwise together with their adjoints.  In particular,
it holds for matrix-valued symbols which are simultaneously unitarily
diagonalizable, with constants independent of the matrix size.
\end{corollary}

\section{Proof}

We first record the scalar input in the precise form needed here.

\begin{lemma}[Scalar paraproduct--Carleson equivalence]\label{lem:scalar}
For a scalar $d$-adic martingale $\beta$, let
\[
 C(\beta)^2=\sup_n\left\|\E_n\sum_{k>n}|d_k\beta|^2
                    \right\|_\infty.
\]
Then the scalar paraproduct $\pi_\beta$ on $L_2(\mathbb R)$ is bounded if
and only if $C(\beta)<\infty$, and
\[
 C(\beta)\le \|\pi_\beta\|_{2\to2}\le 2C(\beta).
\]
Furthermore,
\[
 C(\beta)\asymp_d
 \sup_n\left\|\E_n\sum_{k\ge n}|d_k\beta|^2\right\|_\infty^{1/2}.
\]
The same estimates hold for Hilbert-space-valued inputs, since the scalar
paraproduct is tensored with the identity.
\end{lemma}

\begin{proof}
The quantity $C(\beta)^2$ is the scalar martingale Carleson constant of the
positive sequence $a_k=|d_k\beta|^2$.  The usual stopping-time proof of the
Carleson embedding theorem gives
\[
 \sum_k\int a_k\,|\E_{k-1}f|^2\,dm
 \le 4C(\beta)^2\|f\|_2^2.
\]
The left side equals $\|\pi_\beta f\|_2^2$, because the summands
$d_k\beta\,\E_{k-1}f$ are orthogonal martingale differences.  This proves the
upper estimate.  Conversely, test on indicators of finite-measure
$d$-adic sets at level $n$.  For all later levels the conditional expectation
of such an indicator is the indicator itself, so boundedness of $\pi_\beta$
gives the Carleson inequality and hence
$C(\beta)\le\|\pi_\beta\|$.

It remains only to compare the two indexing conventions.  The $k>n$ sum is
contained in the $k\ge n$ sum.  In the other direction, on every parent
$d$-adic atom, $d_n\beta$ is constant on each of the $d$ children and has mean
zero.  Therefore its largest child value squared is at most $d$ times its
parent average:
\[
 |d_n\beta|^2\le d\,\E_{n-1}|d_n\beta|^2.
\]
Taking the appropriate conditional sums yields the asserted comparison with
a constant depending only on $d$.  Tensoring a bounded scalar operator with
an identity preserves its norm, which proves the Hilbert-valued statement.
\end{proof}

\begin{proof}[Proof of Theorem~\ref{thm:main}]
Put $H=L_2(\M,\tau)$, and represent $\A$ faithfully on $H$ by left
multiplication $a\mapsto L_a$.  This is an abelian von Neumann algebra of
normal operators.  By the spectral representation theorem for abelian von
Neumann algebras \cite[Ch.~9]{KadisonRingrose}, there are a measure space
$(\Omega,\nu)$, a measurable field of Hilbert spaces $(H_\omega)$, and a
unitary
\[
 U:H\longrightarrow\int_\Omega^\oplus H_\omega\,d\nu(\omega)
\]
such that every $L_a$, $a\in\A$, acts in the fibers as multiplication by a
scalar function $\widehat a(\omega)$.

Apply this representation to the strongly measurable symbol $b$.  After
discarding a null set, write
\[
 U L_{b(t)}U^{-1}\big|_{H_\omega}
   =\beta_\omega(t) I_{H_\omega}.
\]
Time conditional expectations commute with this fixed fiber decomposition;
thus $d_kb$ corresponds to $d_k\beta_\omega$.

For integers $r\le s$, let
\[
 \pi_b^{[r,s]}f=\sum_{k=r}^s d_kb\,\E_{k-1}f.
\]
Conjugating by $I_{L_2(\mathbb R)}\otimes U$ gives the decomposable operator
\[
 \int_\Omega^\oplus
   \bigl(\pi_{\beta_\omega}^{[r,s]}\otimes I_{H_\omega}\bigr)
 \,d\nu(\omega).
\]
Consequently,
\begin{equation}\label{eq:fibernorm}
 \|\pi_b^{[r,s]}\|
 =\operatorname*{ess\,sup}_{\omega\in\Omega}
   \|\pi_{\beta_\omega}^{[r,s]}\|.
\end{equation}

The positive operator
$\E_n\sum_{k\ge n}|d_kb|^2$ lies in $\A$ and has fiber value
$\E_n\sum_{k\ge n}|d_k\beta_\omega|^2$.  The norm of an element of an
abelian von Neumann algebra is the essential supremum of its fiber values.
Since the levels are countable, the two suprema can be interchanged, giving
\begin{equation}\label{eq:fiberbmo}
 \|b\|_{\BMO_{c,d}}
 =\operatorname*{ess\,sup}_{\omega\in\Omega}
   \|\beta_\omega\|_{\BMO_d}.
\end{equation}

If the right side of \eqref{eq:fiberbmo} is finite, Lemma~\ref{lem:scalar}
and \eqref{eq:fibernorm} give a bound, independent of $r,s$,
\[
 \|\pi_b^{[r,s]}\|\le C_d\|b\|_{\BMO_{c,d}}.
\]
On finite martingales the truncations stabilize.  Uniform boundedness
therefore defines a bounded operator $\pi_b$ on the dense subspace of finite
martingales, and then on all of $L_2$, proving the upper estimate in
\eqref{eq:equiv}.

Conversely, if $\pi_b$ is bounded, every finite interval truncation is bounded
by $\|\pi_b\|$ (equivalently, one may project its output onto the corresponding
orthogonal martingale-difference levels).  Equations
\eqref{eq:fibernorm}--\eqref{eq:fiberbmo}, the lower scalar estimate, and
monotone passage through finite level intervals yield
\[
 \|b\|_{\BMO_{c,d}}\le c_d^{-1}\|\pi_b\|.
\]
This proves the theorem.
\end{proof}

\section{Verification notes and limitations}

\begin{itemize}
\item The proof is dimension-free because diagonalization produces scalar
fibers before the scalar Carleson estimate is applied.
\item The ambient algebra $\M$ need not be abelian.  Inputs remain
$L_2(\M)$-valued; the abelian hypothesis concerns only the von Neumann algebra
generated by the symbol range.
\item The proof uses finite level truncations, so no unjustified pointwise
convergence of the infinite paraproduct series is required.
\item The theorem does not address noncommuting symbols.  In that regime the
fiber reduction fails, and the logarithmic matrix obstruction cited by Wei
and Zhang remains relevant.
\item The most useful human check is the measurable direct-integral step for
the desired level of generality (especially nonseparable ambient algebras).
For the standard separable-predual setting, or for essentially separably
valued symbols, the displayed decomposition is routine.  The same proof can
be phrased inside the abelian von Neumann tensor product to avoid separability.
\end{itemize}

\paragraph{Novelty check.}
The run's lightweight indexes and bounded arXiv searches performed on
19 July 2026 for the phrases ``operator-valued martingale paraproduct commuting
symbol'', ``semicommutative paraproduct abelian'', and ``commuting normal
operator symbol BMO'' found no explicit statement of Theorem~\ref{thm:main}.
The result may be folklore implicit in the abelian spectral theorem and scalar
Carleson embedding, so this packet claims a likely-valid partial result, not a
high-confidence literature novelty claim.

\begin{thebibliography}{9}
\bibitem{WeiZhang}
Z.~Wei and H.~Zhang,
\emph{Boundedness of operator-valued commutators involving martingale
paraproducts}, arXiv:2401.08729 (2024).

\bibitem{KadisonRingrose}
R.~V. Kadison and J.~R. Ringrose,
\emph{Fundamentals of the Theory of Operator Algebras, Vol.~II},
Academic Press, 1986.

\bibitem{NPTV}
F.~Nazarov, G.~Pisier, S.~Treil, and A.~Volberg,
\emph{Sharp estimates in vector Carleson imbedding theorem and for vector
paraproducts}, J. Reine Angew. Math. 542 (2002), 147--171.
\end{thebibliography}

\end{document}
